% ════════════════════════════════════════════════════════════════
\subsection{Egzamin zerowy}
% ════════════════════════════════════════════════════════════════

\begin{zadanie}[Zadanie 1 --- dwa najtańsze symbole pod wspólnym rodzicem (15 p.)]
	Udowodnij, że w~kodowaniu Huffmana istnieje optymalne drzewo kodowe,
	w~którym dwa najtańsze (najrzadziej występujące) symbole mają wspólnego
	rodzica.
\end{zadanie}

\begin{zadanie}[Zadanie 2 --- skojarzenie najliczniejsze a~ścieżki powiększające (15 p.)]
	Udowodnij, że skojarzenie w~grafie jest najliczniejsze wtedy i~tylko wtedy,
	gdy nie istnieje ścieżka powiększająca.
\end{zadanie}

\begin{zadanie}[Zadanie 3 --- matroid jednego wierzchołka na składową (15 p.)]
	Niech $G$ będzie grafem, a~$\mathcal{A}_G$ rodziną podzbiorów $V(G)$ taką, że
	$S\in\mathcal{A}_G$ wtedy i~tylko wtedy, gdy każda składowa spójności $G$
	zawiera co najwyżej jeden wierzchołek z~$S$. Rozstrzygnij, czy para
	$(V(G),\mathcal{A}_G)$ zawsze jest matroidem (udowodnij to lub podaj
	kontrprzykład).
\end{zadanie}

\begin{zadanie}[Zadanie 4 --- dokładne pokrycie wierzchołkowe w~czasie podwykładniczym (15 p.)]
	Zaprojektuj algorytm dokładny, który znajdzie liczność najmniejszego pokrycia
	wierzchołkowego grafu w~czasie $o(2^n)$, gdzie $n$ jest liczbą wierzchołków.
	Uzasadnij poprawność i~oszacuj złożoność swojego algorytmu.
\end{zadanie}

\begin{zadanie}[Zadanie 5 --- 3-aproksymacja problemu 3DC (15 p.)]
	Rozważmy problem --- nazwijmy go roboczo 3DC --- w~którym mamy daną rodzinę
	$\mathcal{R}$ 3-elementowych podzbiorów pewnego zbioru $X$ i~poszukujemy jak
	najmniej licznego zbioru $S\subseteq X$ takiego, że każdy zbiór z~$\mathcal{R}$
	zawiera przynajmniej jeden element z~$S$. Zaprojektuj wielomianowy algorytm
	$3$-aproksymacyjny dla problemu 3DC (i~udowodnij, że stała aproksymacji
	istotnie jest nie większa niż~$3$). Możesz zainspirować się algorytmem
	$2$-aproksymacyjnym dla problemu najmniejszego pokrycia wierzchołkowego.
\end{zadanie}
